% Section 1.3, from lectures/L01/appendix/c_counting_cycles.md.

\section{Counting Cycles}
\label{c1:sec:cycles}

\subsection{Why bother}
\label{c1:sec:whybother}
On most machines you cannot say what a routine costs. Caches, branch prediction and out-of-order
execution mean the same code takes different times on different runs, and the honest answer to
``how long does this take'' is ``measure it, repeatedly, and take a distribution''.

The ATmega328P has none of those. Every instruction takes a fixed number of cycles, the number is
in a table, and it is the same number every time. So ``how long does this take'' has an exact
answer you can work out on paper before the code exists, and that changes what is worth doing:
\begin{itemize}
  \item You can decide between two implementations without writing either.
  \item You can tell whether a delay loop will meet its deadline without a scope.
  \item When a measurement disagrees with your count, one of you is wrong about the code, and
    finding out which is a productive afternoon rather than a guess about the hardware.
\end{itemize}

This section is that table, and the arithmetic on top of it.

\subsection{What each instruction costs}
\label{c1:sec:costs}
Every figure here comes from the AVR Instruction Set Manual for this core, and every one was
confirmed by running that instruction in the simulator and subtracting the \code{ret} around it.
Where the two had disagreed, the disagreement would be written down here. They did not.

\begin{booktable}{@{}p{3.6em}L@{}}
  \thead{Cycles} & \thead{Instructions} \\
  \midrule
  1 & \code{ldi mov movw add adc sub subi and andi or ori eor com neg cp cpc cpi lsl lsr rol ror
      asr inc dec nop in out sei cli} \\
  2 & \code{rjmp lds sts ld ldd st std push pop adiw sbiw sbi cbi} \\
  3 & \code{rcall lpm} \\
  4 & \code{ret reti} \\
  1 or 2 & \code{breq brne brlo brsh brcc brcs} and the other conditional branches \\
\end{booktable}

Three rows in that table are worth more attention than the rest.

\textbf{\code{ret} costs four cycles and \code{rcall} costs three.} Calling a subroutine that does
nothing at all costs seven cycles, which is more than most loops' bodies. That is not an argument
against subroutines; it is the reason a driver's cost is never just the cost of its body, and the
reason the cross-check in Exercise~\ref{c1:ex:crosscheck} comes out the way it does.

\textbf{\code{sbi} and \code{cbi} cost two cycles, not one.} They look like single-bit versions of
\code{out}, and they are not: they read the register, modify one bit, and write it back. Worth
remembering when you are counting a routine that sets several bits one at a time and wondering
where the cycles went.

\textbf{A memory access costs twice what a register operation does.} \code{lds} is two cycles and
\code{mov} is one, which is the whole argument for keeping working values in registers. On a
machine with 32 of them that is usually possible, and arranging for it is most of what makes
assembly fast.

\subsection{A branch costs more when it branches}
\label{c1:sec:branchcost}
A conditional branch takes \textbf{one cycle if it falls through and two if it is taken}. The
reason is that a taken branch changes the program counter, and the instruction already being
fetched has to be thrown away.

This is the single most common source of an off-by-a-few cycle count, because it means a loop's
cost is not the number of iterations times the body.

Take a loop with its test at the top, which is the shape the subroutine in \secref{c1:sec:first}
needs: compare a running count against the argument, branch out of the loop when they are equal,
do the work, add one to the count, and jump back to the top.

\begin{booktable}{@{}Lp{10.6em}l@{}}
  \thead{Step} & \thead{Cycles} & \thead{Times it runs} \\
  \midrule
  compare the count against the argument & 1 & $n + 1$ \\
  branch out when they are equal & 1 falling through, 2 taken & $n + 1$ \\
  the work, one instruction here & 1 & $n$ \\
  add one to the count & 1 & $n$ \\
  jump back to the top & 2 & $n$ \\
\end{booktable}

A loop running $n$ times executes the branch \textbf{$n + 1$ times}: $n$ of them fall
through into the body, and one is taken to leave. So the branch contributes
$n \times 1 + 1 \times 2 = n + 2$ cycles, not $n + 1$ and not $2n$.

Everything else in the loop runs exactly $n$ times, except the comparison, which also runs on the
final pass that leaves the loop and so runs $n + 1$ times as well. That extra comparison is the one
people forget, and it is why the answer comes out ten rather than nine when $n$ is zero.

\subsection{Counting a routine on paper}
\label{c1:sec:onpaper}
The table in \secref{c1:sec:branchcost} is the method, and it is worth stating as a method because
you will use it in every chapter of this book.

\textbf{One row per instruction, not per line executed.} Write the instruction, its cost from
\secref{c1:sec:costs}, and how many times it runs. Multiply and sum. A branch gets two rows, because
it costs differently on the passes that fall through and on the one that is taken.

For \code{shift_bits} with an argument of $n$:

\begin{narrowtable}{@{}lccc@{}}
  \thead{Instruction} & \thead{Cycles} & \thead{Times} & \thead{Total} \\
  \midrule
  \code{ldi} $\times$ 2 & 1 & 2 & 2 \\
  \code{cp} & 1 & $n + 1$ & $n + 1$ \\
  \code{breq}, falling through & 1 & $n$ & $n$ \\
  \code{breq}, taken & 2 & 1 & 2 \\
  \code{lsl} & 1 & $n$ & $n$ \\
  \code{inc} & 1 & $n$ & $n$ \\
  \code{rjmp} & 2 & $n$ & $2n$ \\
  \code{mov} & 1 & 1 & 1 \\
  \code{ret} & 4 & 1 & 4 \\
\end{narrowtable}

\noindent which adds to $6n + 10$.

\textbf{Order does not matter and that is not a simplification.} Cost is a sum, so a table that
tracked which instruction ran when would be carrying weight for nothing. What the method does
require is that you have already worked out \textbf{how many times each line runs}, and that is
the part that is actually hard. \Secref{c1:sec:branchcost} is about exactly that, and the $n + 1$ on
the comparison is where most wrong answers come from.

\subsection{What a hand count is blind to}
\label{c1:sec:blind}
It is worth being explicit, because a method that quietly stops being right is worse than no
method.

\textbf{It does not know about the \code{rcall} that reached the routine.} Your table is the body.
The three cycles of getting there, and the instructions that set up the arguments, belong to the
caller, and counting them is the caller's table's job.

\textbf{It does not know about interrupts.} An interrupt arriving mid-routine adds the handler's
cost plus six to nine cycles of entry and four of \code{reti}, at a moment nothing in this model can
predict. From Chapter~3 onwards that is a real effect and a hand count will simply be wrong about
it, by an amount that depends on what else is running. Where that matters, the answer is to
measure.

\textbf{It does not know whether your table matches your code.} Nothing checks that the rows you
wrote describe the routine you assembled. Two things guard against that: the simulator, which
costs the real thing, and the discipline of writing the table from the source rather than from
memory. When they disagree, the table is usually the one that is wrong, because it was written by
a human who was thinking about something else.

\textbf{It says nothing about correctness.} A routine can cost exactly what you predicted and
return entirely the wrong answer. Cycle counting tells you whether something fits in the time
available, and nothing whatever about whether it should be running at all.
